\subsection{Follow the tax difference into earnings}
\label{sec:compound-tax-rollforward}

The initial charge to equity does not tell us the later cost of financing.
That cost depends on how the debt grows and what the tax authority allows the
issuer to deduct. To expose the distinction, consider a separate July issue
whose principal at October maturity is 100, with a coupon of 4. Its proceeds
are again 100 and its initial debt value is 92. The supplied comparable rate
for the 92-day period is now $3/23$, so
\[
 L_0=\frac{100+4}{1+3/23}=92,\qquad
 I=92\frac{3}{23}=12,\qquad L_{\mathrm{after\ coupon}}=92+12-4=100.
\]
The accounting interest of 12 includes both the coupon of 4 and the debt
discount of 8. These are different terms from the preceding principal-92
instrument. At maturity the holder may instead convert the principal into
six ordinary shares after receiving the coupon; the original conversion
equity remains equity under IAS~32.AG32.

Now supply the fiscal treatment. The tax authority treats the initial 100
entirely as debt and permits a deduction for the coupon when it accrues.
It allows no separate deduction for the accounting discount. The fiscal
liability therefore rises to 104 before payment and falls to 100 when the
coupon is paid. Assume sufficient taxable income to use the deduction and
no additional tax on either maturity redemption or original-terms conversion.
These are explicit illustrative fiscal facts, not consequences of a known
tax rate. At the supplied rate of 25\%, they give the following sequence.

\begin{center}
\begin{tabular}{lrrr}
\toprule
Stage & Debt carrying amount & Tax base & Deferred tax liability\\
\midrule
Issue & 92 & 100 & 2\\
Interest accrued & 104 & 104 & 0\\
Coupon paid & 100 & 100 & 0\\
Principal settled & 0 & 0 & 0\\
\bottomrule
\end{tabular}
\end{center}

The first row reduces conversion equity from 8 to 6. The next row has a
different earnings consequence: debit deferred tax liability 2 and credit
deferred tax income 2. IAS~12.23 explicitly places these subsequent changes
in profit or loss. The coupon deduction also reduces current tax expense
by $4/4=1$. Thus the financing cost after tax is
\[
 I+\Delta T_{\mathrm{current}}+\Delta T_{\mathrm{deferred}}
 =12-1-2=9.
\]
There is no second expense of 2 for the initial equity charge. Conversely,
using only the coupon would miss the discount accretion and its tax reversal
\citep[IAS~12, paragraphs 23, 58 and 61A; IAS~33, paragraphs 33--34]{eu20231803}.

The relation extends to the specified one-period fiscal treatment. Let $P$
denote both proceeds and principal at maturity, $L$ the initial debt, $C$ the
coupon and $t$ the unchanged tax rate. Gross conversion equity is $E=P-L$;
effective interest is $I=P+C-L=E+C$. The initial deferred liability $tE$
reverses as book and tax amounts meet, while the coupon supplies a current
deduction of $C$. Therefore
\[
 I-tC-tE=(1-t)I.
\]
A zero coupon still leaves discount accretion and deferred tax income. If
$E=0$, the deferred difference vanishes and only the coupon deduction remains.
Different fiscal bases, changing rates or restricted deductions require their
own schedule before this identity can be used.

Return to the same parent and subsidiary, preserving all their original
operations and ownership changes. The new finance expense reduces group
pretax profit from 136 to 124. Current tax falls from 28 to 27 and deferred
tax expense from 6 to 4, leaving group profit $124-27-4=93$.
Non-controlling interests still receive 7.35 because this financing belongs
to the parent, so parent ordinary earnings are $93-7.35=85.65$.
Conversion leaves cash $644+100-4+1=741$ and equity
$724+6+100-9=821$. Cash redemption instead leaves cash
$644+100-4-100+1=641$ and equity $724+6-9=721$.
The issue-date tax charge is visible in the net equity contribution of 6;
the later reversal is already included in the earnings cost of 9.

Diluted EPS uses the consequential earnings change as well as the dated
shares. Under the explicit assumption of no other earnings or tax effects,
assumed conversion removes interest 12, current tax relief 1 and deferred
tax income 2. Its earnings addback is consequently 9, restoring parent
earnings to $85.65+9=94.65$. The initial equity charge is outside that
numerator. Six potential shares still run for 92 days; actual shares after
October appear only in the conversion world. Using the same day counts as
before gives
\[
 E_{\mathrm{dil,conv}}=\frac{94.65}{4754/365}
   =\frac{138189}{19016},\qquad
 E_{\mathrm{dil,redeem}}=\frac{94.65}{4202/365}
   =\frac{138189}{16808}.
\]
The incremental ratio $9/(552/365)=1095/184$ is below both worlds' basic
EPS, so the series is dilutive under continuing-operation control. The
diluted fractions happen to equal those of the earlier instrument because
both addbacks restore the same earnings and use the same share dates.
The contracts have different actual earnings and principal obligations;
equal diluted EPS does not make them economically interchangeable
\citep[IAS~33, paragraphs 32--38 and 41]{eu20231803}.

The documentary \code{CompoundTaxRollforward/1} calculation reports separate
carrying amounts, fiscal bases, tax origins and journal movements.
\code{CompanyTaxEPS/1} resolves the fixed company inputs and reconstructs the
legal books before accepting an EPS projection. It refuses a coupon-only
addback, a missing deferred-tax consequence, the wrong parent earnings or
overlapping share periods. Missing fiscal findings remain unresolved, and
different rights or tax schedules require another supported method. These
calculations specify the supplied one-period case; early-conversion measurement,
changing tax rates, loss limitations and operative jurisdictional authority
remain separate obligations.
